% ==========================================================================
% INFO 4617 Web Data Science -- Week 14: Automating Data Collection (Chapter 14)
% Build:  latexmk -pdf week-14.tex   (from this folder; no env vars needed)
% (or run `make` from the slides/ directory to build every deck)
% ==========================================================================
\documentclass[9pt,aspectratio=169]{beamer}
% Resolve the shared theme/preamble in ../common/ when compiling this
% file directly from its own folder (no TEXINPUTS needed).
\makeatletter\def\input@path{{../common/}}\makeatother
\input{preamble}
% Larger code than the house default: a frame holds one short block (AUTHORING.md, "Code style")
\lstset{basicstyle=\ttfamily\footnotesize}
% A step of the week's pipeline, for the at-a-glance diagram
\tikzset{stage/.style={draw=black!70,rounded corners=2pt,align=center,inner sep=4pt,font=\footnotesize,minimum height=2.6em,text width=6.3em}}

\title[Week 14 · Automation]{{\Huge Automating Data Collection}}
\subtitle{Week 14 · Practice module · Textbook Chapter 14}
\author[Keegan]{\textbf{Brian C. Keegan, Ph.D.} \\ Associate Professor, Department of Information Science \\ University of Colorado Boulder}
\date{November 18 \& 20, 2026}

\begin{document}

\maketitle

% -------------------------------------------------------------------------
% Schedule: Monday, Nov 16 finishes week 13's lab (Part 3, in week 13's deck). Week 14's lab runs
% Wednesday, Nov 18 (Part 1) and Friday, Nov 20 (Part 2, then 15 minutes on portfolio checkpoint 4).
% The textbook's plan for chapter 14: docs/plans/2026-11-13-ch14-automation.md.
%
% Code frames follow the lab key, labs/week-14/lab-14-key.ipynb, and the template it prints,
% labs/week-14/template/ (collector.py and .github/workflows/collect.yml), as written on
% 2026-11-13 (week 14's lab branch), reflowed to fit and sometimes cut with "...". Some of the
% template's own comments are shortened into the listings' one-line comments. The two my_series
% frames show the address a student types (the cell's comment), where the key reads the course's
% series. No frame shows printed output: the slides lead into the notebook, where students run each
% cell (instructor, 2026-11-12).
% One code block a frame, narration on the left and code on the right: slides/common/AUTHORING.md.
% Code shaped for reading (2026-11-13): a one-line comment above each step, at most two chained
% methods, one argument per line in long calls. Where a slide names a step the key or the template
% doesn't, it does the same thing and prints the same things; its narration says so.
\begin{frame}{This week at a glance}
    \begin{center}
    \begin{tikzpicture}[node distance=0.35cm, >=stealth]
        \node[stage] (one) {\textbf{One request}\\[-1pt]{\scriptsize by hand · Wed}};
        \node[stage, right=of one] (script) {\textbf{collector.py}\\[-1pt]{\scriptsize a script · Wed}};
        \node[stage, right=of script] (actions) {\textbf{Every hour}\\[-1pt]{\scriptsize Actions · Wed}};
        \node[stage, right=of actions] (csv) {\textbf{forecasts.csv}\\[-1pt]{\scriptsize read it · Fri}};
        \node[stage, right=of csv] (fail) {\textbf{A red X}\\[-1pt]{\scriptsize break it · Fri}};
        \draw[->, thick] (one) -- (script);
        \draw[->, thick] (script) -- (actions);
        \draw[->, thick] (actions) -- (csv);
        \draw[->, thick] (csv) -- (fail);
    \end{tikzpicture}
    \end{center}
    \begin{columns}[T]
        \begin{column}{0.6\textwidth}
            \dayband{Monday, Nov 16}{Week 13's lab, Part 3}\par
            {\small Its frames are in week 13's deck.\par}
            \medskip
            \dayband{Wednesday, Nov 18}{Part 1}\par
            {\small One request by hand, then your own collector running in GitHub Actions.\par}
            \medskip
            \dayband{Friday, Nov 20}{Part 2}\par
            {\small What it gathered, failures and monitoring. Then 15 minutes on portfolio checkpoint 4.\par}
        \end{column}
        \begin{column}{0.35\textwidth}
            \textbf{Bring}\\[2pt]
            {\footnotesize Your GitHub account, and \texttt{boulder-forecast-template.zip} from Canvas. No key.\par}
            \medskip
            \textbf{Due}\\[2pt]
            {\footnotesize Sunday, Nov 22: chapter 14's Recommended Exercises.\par}
            {\footnotesize Sunday, Nov 29: portfolio checkpoint 4.\par}
        \end{column}
    \end{columns}
\end{frame}

% =========================================================================
\section{Wednesday · Live-Coding Lab}
% =========================================================================

% Placeholder: fill in from the daily-note responses before class
% (python3 ../common/check_daily_questions.py week-14).

\begin{frame}{One request is a snapshot}
    \begin{columns}[T]
        \begin{column}{0.6\textwidth}
            Forecasters change Boulder's forecast a few times a day. On November 12, one forecast stood for 12 hours. One request tells you what it said once.

            \medskip
            To see it change, a program has to ask again and again, on a schedule.

            \medskip
            What else from this semester would you want to watch change?
        \end{column}
        \begin{column}{0.35\textwidth}
            \textbf{This week's source}\\[2pt]
            {\footnotesize The National Weather Service's hourly forecast for Boulder. No key.\par}
        \end{column}
    \end{columns}
\end{frame}

\begin{frame}{This week's lab}
    \begin{columns}[T]
        \begin{column}{0.6\textwidth}
            \dayband{Find it}{\texttt{labs/week-14} in the course repository}
            \begin{enumerate}\small
                \item Open \texttt{lab-14.ipynb} and click \textbf{Download raw file}.
                \item Open it in JupyterLab, in the \texttt{webdata} environment.
                \item Sign in to GitHub in your browser.
            \end{enumerate}
            \medskip
            \dayband{Use it}{one part a day}
            \begin{itemize}\small
                \item I type the code. You type it in your copy.
            \end{itemize}
        \end{column}
        \begin{column}{0.35\textwidth}
            \textbf{Fall behind?}\\[2pt]
            {\footnotesize Copy the cell from \texttt{lab-14-key.ipynb}, beside it, and keep going.\par}
        \end{column}
    \end{columns}
\end{frame}

\begin{frame}[fragile]{Step 0: set up}
    \begin{columns}[T]
        \begin{column}{0.37\textwidth}
            This cell is given. Put your email in \texttt{HEADERS}, and run it.

            \medskip
            The NWS asks for no key. It asks every program to name itself and give a contact. Why send your email if nobody checks?

            \medskip
            {\scriptsize\textbf{Docs:} the NWS's \href{https://www.weather.gov/documentation/services-web-api}{API documentation}, Overview tab\par}
        \end{column}
        \begin{column}{0.6\textwidth}
\begin{lstlisting}
# The two libraries we use
import requests
import pandas as pd
# Name our program, and give your email
HEADERS = {"User-Agent":
           "WebDataScience/1.0 (your-email@colorado.edu)"}
\end{lstlisting}
        \end{column}
    \end{columns}
\end{frame}

\begin{frame}[fragile]{Step 0: two addresses}
    \begin{columns}[T]
        \begin{column}{0.37\textwidth}
            The same cell names two addresses. \texttt{nws\_base} is the NWS's API. \texttt{template} is where the course keeps \texttt{collector.py} and its workflow.

            \medskip
            Steps 3 and 4 print those two files from it.
        \end{column}
        \begin{column}{0.6\textwidth}
\begin{lstlisting}
# The NWS's API
nws_base = "https://api.weather.gov"
# Where the course keeps this week's template
template = (
    "https://raw.githubusercontent.com/"
    "cuinfoscience/INFO4617-Fall2026/main/"
    "labs/week-14/template/"
)
\end{lstlisting}
        \end{column}
    \end{columns}
\end{frame}

\begin{frame}[fragile]{Step 1: where is Boulder?}
    \begin{columns}[T]
        \begin{column}{0.37\textwidth}
            The NWS forecasts for a grid of squares about 2.5 km on a side.

            \medskip
            \texttt{/points} takes a latitude, then a longitude, and names the square. West is negative. Which number is the latitude?
        \end{column}
        \begin{column}{0.6\textwidth}
\begin{lstlisting}
# Boulder: the latitude, then the longitude
points_url = f"{nws_base}/points/40.015,-105.2705"
\end{lstlisting}
        \end{column}
    \end{columns}
\end{frame}

\begin{frame}[fragile]{Step 1: send it}
    \begin{columns}[T]
        \begin{column}{0.37\textwidth}
            Send it with \texttt{HEADERS}. Print the status code and the \texttt{Content-Type} header.

            \medskip
            Is the answer JSON? Which kind?
        \end{column}
        \begin{column}{0.6\textwidth}
\begin{lstlisting}
# Send it, with our headers
points_response = requests.get(
    points_url,
    headers=HEADERS,
)
# The status code, and the kind of data
print(
    points_response.status_code,
    points_response.headers["Content-Type"],
)
\end{lstlisting}
        \end{column}
    \end{columns}
\end{frame}

\begin{frame}[fragile]{Step 1: any error?}
    \begin{columns}[T]
        \begin{column}{0.37\textwidth}
            When something goes wrong, the NWS says what in \texttt{detail}. \texttt{None} means no error.

            \medskip
            Bluesky's errors had a \texttt{message}. Why read the error before the data?
        \end{column}
        \begin{column}{0.6\textwidth}
\begin{lstlisting}
# The answer, as a dictionary
points = points_response.json()
# Any error? None means no error
print(points.get("detail"))
\end{lstlisting}
        \end{column}
    \end{columns}
\end{frame}

\begin{frame}[fragile]{Step 1: one level at a time}
    \begin{columns}[T]
        \begin{column}{0.37\textwidth}
            Print the answer's keys. Then print the keys of its \texttt{properties}.

            \medskip
            Which key looks like the address of an hourly forecast?
        \end{column}
        \begin{column}{0.6\textwidth}
\begin{lstlisting}
# The answer's keys
print(points.keys())
# The keys one level down, in "properties"
print(points["properties"].keys())
\end{lstlisting}
        \end{column}
    \end{columns}
\end{frame}

\begin{frame}[fragile]{Step 1: the square}
    \begin{columns}[T]
        \begin{column}{0.37\textwidth}
            Each square belongs to one of the NWS's forecast offices.

            \medskip
            Which office, which square, and which city did you get?
        \end{column}
        \begin{column}{0.6\textwidth}
\begin{lstlisting}
# The part of the answer about the square
grid = points["properties"]
# Its office, and the square's two numbers
print(
    grid["gridId"],
    grid["gridX"],
    grid["gridY"],
)
# The city nearest the square
print(grid["relativeLocation"]["properties"]["city"])
\end{lstlisting}
        \end{column}
    \end{columns}
\end{frame}

\begin{frame}[fragile]{Step 1: the forecast's address}
    \begin{columns}[T]
        \begin{column}{0.37\textwidth}
            \texttt{forecastHourly} is the address our collector will ask every hour.

            \medskip
            The documentation says a square can change. \texttt{collector.py} keeps only this address. What would it collect if the square changed?
        \end{column}
        \begin{column}{0.6\textwidth}
\begin{lstlisting}
# The address of the hourly forecast
hourly_url = grid["forecastHourly"]
print(hourly_url)
\end{lstlisting}
        \end{column}
    \end{columns}
\end{frame}

\begin{frame}[fragile]{Microsprint: the numbers swapped}
    \begin{columns}[T]
        \begin{column}{0.37\textwidth}
            The likeliest slip is the order of the two numbers. Ask with the longitude first.

            \medskip
            What status code came back? What did the NWS say was wrong, in its own words?
        \end{column}
        \begin{column}{0.6\textwidth}
\begin{lstlisting}
# The longitude first, by mistake
swapped_response = requests.get(
    f"{nws_base}/points/-105.2705,40.015",
    headers=HEADERS,
)
# The answer, as a dictionary
swapped = swapped_response.json()
# The status code, and the NWS's error
print(swapped_response.status_code, swapped.get("detail"))
# Which part of the request was wrong
print(swapped.get("parameterErrors"))
\end{lstlisting}
        \end{column}
    \end{columns}
\end{frame}

\begin{frame}[fragile]{Step 2: the hourly forecast}
    \begin{columns}[T]
        \begin{column}{0.37\textwidth}
            Ask \texttt{hourly\_url}, the same way.

            \medskip
            Is the status code the same as for \texttt{/points}?
        \end{column}
        \begin{column}{0.6\textwidth}
\begin{lstlisting}
# Send it, the same way
hourly_response = requests.get(
    hourly_url,
    headers=HEADERS,
)
# The status code, and the kind of data
print(
    hourly_response.status_code,
    hourly_response.headers["Content-Type"],
)
\end{lstlisting}
        \end{column}
    \end{columns}
\end{frame}

\begin{frame}[fragile]{Step 2: one period per hour}
    \begin{columns}[T]
        \begin{column}{0.37\textwidth}
            Its \texttt{properties} hold a list of \texttt{periods}, one for each hour ahead.

            \medskip
            How many periods? How many days is that?
        \end{column}
        \begin{column}{0.6\textwidth}
\begin{lstlisting}
# The answer, as a dictionary
hourly_json = hourly_response.json()
# The part that holds the forecast
hourly = hourly_json["properties"]
# How many periods it holds
print(len(hourly["periods"]), "periods")
\end{lstlisting}
        \end{column}
    \end{columns}
\end{frame}

\begin{frame}[fragile]{Step 2: one hour's forecast}
    \begin{columns}[T]
        \begin{column}{0.37\textwidth}
            Print the first period.

            \medskip
            Which of its fields would you save every hour? Which would you leave out?
        \end{column}
        \begin{column}{0.6\textwidth}
\begin{lstlisting}
# The first period: one hour's forecast
print(hourly["periods"][0])
\end{lstlisting}
        \end{column}
    \end{columns}
\end{frame}

\begin{frame}[fragile]{Step 2: how new is it?}
    \begin{columns}[T]
        \begin{column}{0.37\textwidth}
            \texttt{updateTime}: when forecasters last changed the forecast.

            \medskip
            \texttt{generatedAt}: when the NWS's server built this answer.

            \medskip
            Which is older? By how much?
        \end{column}
        \begin{column}{0.6\textwidth}
\begin{lstlisting}
# When forecasters last changed the forecast
print("updated:  ", hourly["updateTime"])
# When the NWS's server built this answer
print("generated:", hourly["generatedAt"])
\end{lstlisting}
        \end{column}
    \end{columns}
\end{frame}

\begin{frame}[fragile]{Step 2: what the headers say}
    \begin{columns}[T]
        \begin{column}{0.37\textwidth}
            \texttt{Last-Modified} dates this answer. \texttt{Cache-Control} says how long a copy may be reused: \texttt{max-age} by your browser, \texttt{s-maxage} by the NWS's caches.

            \medskip
            Thirty of us just asked. How many times did the NWS build it?

            \medskip
            {\scriptsize\textbf{Docs:} MDN's \href{https://developer.mozilla.org/en-US/docs/Web/HTTP/Reference/Headers/Cache-Control}{Cache-Control}\par}
        \end{column}
        \begin{column}{0.6\textwidth}
\begin{lstlisting}
# When this answer was last changed
print(
    "Last-Modified:",
    hourly_response.headers["Last-Modified"],
)
# How long a copy of it may be reused
print(
    "Cache-Control:",
    hourly_response.headers["Cache-Control"],
)
\end{lstlisting}
        \end{column}
    \end{columns}
\end{frame}

\begin{frame}[fragile]{Microsprint: how old is the forecast?}
    \begin{columns}[T]
        \begin{column}{0.37\textwidth}
            \texttt{pd.to\_datetime()} turns each time into one Python can subtract.

            \medskip
            How old was your forecast when the NWS built your answer? And your neighbor's?
        \end{column}
        \begin{column}{0.6\textwidth}
\begin{lstlisting}
# Each time, as one Python can subtract
generated = pd.to_datetime(hourly["generatedAt"])
updated = pd.to_datetime(hourly["updateTime"])
# How old the forecast was when the answer was built
print(generated - updated)
\end{lstlisting}
        \end{column}
    \end{columns}
\end{frame}

\begin{frame}{Step 3: from a notebook to a script}
    \begin{columns}[T]
        \begin{column}{0.6\textwidth}
            \texttt{collector.py} is the cells above, saved as a script that a computer can run without us. A given cell prints it.

            \medskip
            Find our request in it. Then find four things a notebook didn't need.
        \end{column}
        \begin{column}{0.35\textwidth}
            \textbf{The four}\\[2pt]
            {\footnotesize Constants. Append mode. \texttt{raise\_for\_status()}. The \texttt{\_\_main\_\_} guard.\par}
        \end{column}
    \end{columns}
\end{frame}

\begin{frame}[fragile]{collector.py: constants at the top}
    \begin{columns}[T]
        \begin{column}{0.37\textwidth}
            Each setting is written once, at the top, and changed in one place.

            \medskip
            Which line would you change to collect another city's forecast?
        \end{column}
        \begin{column}{0.6\textwidth}
\begin{lstlisting}
# The hourly forecast for Boulder's grid square
URL = (
    "https://api.weather.gov/gridpoints/BOU/54,75/"
    "forecast/hourly"
)
# Name our program, and give your email
HEADERS = {"User-Agent":
           "WebDataScience/1.0 (your-email@colorado.edu)"}
\end{lstlisting}
        \end{column}
    \end{columns}
\end{frame}

\begin{frame}[fragile]{collector.py: the file, its columns, a last day}
    \begin{columns}[T]
        \begin{column}{0.37\textwidth}
            \texttt{FIELDS} names each row's columns, in order. The \texttt{...} stands for four more.

            \medskip
            Which column would you add? From which field of a period?
        \end{column}
        \begin{column}{0.6\textwidth}
\begin{lstlisting}
# Where the rows go
FILE = "forecasts.csv"
# The columns each row has, in order
FIELDS = [
    "collected_at",
    "update_time",
    "start_time",
    "temperature",
    ...  # four more
]
# The last day to collect
LAST_DAY = date(2026, 12, 11)
\end{lstlisting}
        \end{column}
    \end{columns}
\end{frame}

\begin{frame}[fragile]{collector.py: stop after the last day}
    \begin{columns}[T]
        \begin{column}{0.37\textwidth}
            After December 11, \texttt{main()} prints a line and stops before it asks.

            \medskip
            Why would a collector need a last day? Who would a forgotten one keep asking?
        \end{column}
        \begin{column}{0.6\textwidth}
\begin{lstlisting}
def main():
    # After the last day, stop before asking
    if datetime.now(timezone.utc).date() > LAST_DAY:
        print("Collection ended on", LAST_DAY)
        return
    ...
\end{lstlisting}
        \end{column}
    \end{columns}
\end{frame}

\begin{frame}[fragile]{collector.py: fail where someone will see}
    \begin{columns}[T]
        \begin{column}{0.37\textwidth}
            \texttt{raise\_for\_status()} stops the script with an error if the answer isn't 200.

            \medskip
            Our notebook cells never use it. Why does a script that runs alone need it?

            \medskip
            {\scriptsize\textbf{Docs:} \href{https://requests.readthedocs.io/en/stable/api/\#requests.Response.raise_for_status}{raise\_for\_status()}\par}
        \end{column}
        \begin{column}{0.6\textwidth}
\begin{lstlisting}
    # Ask, and give up after 30 seconds
    response = requests.get(
        URL,
        headers=HEADERS,
        timeout=30,
    )
    # Stop with an error if the answer isn't 200
    response.raise_for_status()
    # The part that holds the forecast
    forecast = response.json()["properties"]
\end{lstlisting}
        \end{column}
    \end{columns}
\end{frame}

\begin{frame}[fragile]{collector.py: when we asked}
    \begin{columns}[T]
        \begin{column}{0.37\textwidth}
            Every row will repeat when we asked and when the NWS last updated the forecast.

            \medskip
            Why keep those two on every row?

            \medskip
            \texttt{collector.py} gets \texttt{collected\_at} in one line.
        \end{column}
        \begin{column}{0.6\textwidth}
\begin{lstlisting}
    # The time now, in UTC
    now = datetime.now(timezone.utc)
    # As text, to the second
    collected_at = now.isoformat(timespec="seconds")
    # When the NWS last updated the forecast
    update_time = forecast["updateTime"]
\end{lstlisting}
        \end{column}
    \end{columns}
\end{frame}

\begin{frame}[fragile]{collector.py: append}
    \begin{columns}[T]
        \begin{column}{0.37\textwidth}
            \texttt{"a"} opens the file to add rows at its end.

            \medskip
            What would \texttt{"w"} do to last week's rows?
        \end{column}
        \begin{column}{0.6\textwidth}
\begin{lstlisting}
    # Is this the first run? Then there's no file yet
    new_file = not os.path.exists(FILE)
    # Open the file to add rows at its end
    with open(
        FILE,
        "a",
        newline="",
        encoding="utf-8",
    ) as f:
        ...  # write the rows: the next two frames
\end{lstlisting}
        \end{column}
    \end{columns}
\end{frame}

\begin{frame}[fragile]{collector.py: write the header once}
    \begin{columns}[T]
        \begin{column}{0.37\textwidth}
            The header goes in only when the file is new.

            \medskip
            \texttt{new\_file} is set before \texttt{open()}. Why not after?

            \medskip
            {\scriptsize\textbf{Docs:} \href{https://docs.python.org/3/library/csv.html\#csv.DictWriter}{csv.DictWriter}\par}
        \end{column}
        \begin{column}{0.6\textwidth}
\begin{lstlisting}
    with open(...) as f:
        # A writer that turns a dictionary into a row
        writer = csv.DictWriter(f, fieldnames=FIELDS)
        # The header, only in a new file
        if new_file:
            writer.writeheader()
        # Then a row for each hour
        for period in forecast["periods"]:
            ...
\end{lstlisting}
        \end{column}
    \end{columns}
\end{frame}

\begin{frame}[fragile]{collector.py: one row per hour}
    \begin{columns}[T]
        \begin{column}{0.37\textwidth}
            Each period becomes one row: a dictionary whose keys are the names in \texttt{FIELDS}.

            \medskip
            What would \texttt{writerow()} do with a key that isn't in \texttt{FIELDS}?
        \end{column}
        \begin{column}{0.6\textwidth}
\begin{lstlisting}
        for period in forecast["periods"]:
            # One row per hour, as a dictionary
            writer.writerow({
                "collected_at": collected_at,
                "update_time": update_time,
                "start_time": period["startTime"],
                "temperature": period["temperature"],
                ...  # four more fields
            })
\end{lstlisting}
        \end{column}
    \end{columns}
\end{frame}

\begin{frame}[fragile]{collector.py: run as a script}
    \begin{columns}[T]
        \begin{column}{0.37\textwidth}
            The last two lines run \texttt{main()} only when the file is run as a script, as the workflow does with \texttt{python collector.py}.

            \medskip
            Where else in the file would you look for a bug?

            \medskip
            {\scriptsize\textbf{Docs:} \href{https://docs.python.org/3/library/__main__.html\#idiomatic-usage}{\texttt{\_\_main\_\_}}\par}
        \end{column}
        \begin{column}{0.6\textwidth}
\begin{lstlisting}
# Run main() only when the file is run as a script
if __name__ == "__main__":
    main()
\end{lstlisting}
        \end{column}
    \end{columns}
\end{frame}

\begin{frame}{Step 4: your own repository}
    \begin{columns}[T]
        \begin{column}{0.6\textwidth}
            \begin{enumerate}\small
                \item Unzip \texttt{boulder-forecast-template.zip}. It is in this week's Canvas module.
                \item On GitHub, click \textbf{+} at the top right. Then click \textbf{New repository}.
                \begin{itemize}\small
                    \item Type the name \texttt{boulder-forecast}. Select \textbf{Public}.
                    \item Turn on the README option. Click \textbf{Create repository}.
                \end{itemize}
            \end{enumerate}
        \end{column}
        \begin{column}{0.35\textwidth}
            \textbf{Why a README?}\\[2pt]
            {\footnotesize It gives the new repository its \texttt{main} branch.\par}
        \end{column}
    \end{columns}
\end{frame}

\begin{frame}{Step 4: upload two files}
    \begin{columns}[T]
        \begin{column}{0.6\textwidth}
            \begin{enumerate}\small\setcounter{enumi}{2}
                \item Click \textbf{Add file}, then \textbf{Upload files}. Drag in \texttt{collector.py} and \texttt{README.md}. Click \textbf{Commit changes}.
                \item Click \texttt{collector.py}, then the pencil. In \texttt{HEADERS}, type your email address. Commit the change.
            \end{enumerate}
        \end{column}
        \begin{column}{0.35\textwidth}
            \textbf{No \texttt{forecasts.csv}?}\\[2pt]
            {\footnotesize The first run makes it.\par}
        \end{column}
    \end{columns}
\end{frame}

\begin{frame}{Step 4: create the workflow}
    \begin{columns}[T]
        \begin{column}{0.6\textwidth}
            \begin{enumerate}\small\setcounter{enumi}{4}
                \item Click \textbf{Add file}, then \textbf{Create new file}.
                \begin{itemize}\small
                    \item Type the name \texttt{.github/workflows/collect.yml}. Each \texttt{/} makes a folder.
                    \item Run the given cell under Step 4. Copy all that it prints into the big box.
                    \item Commit the change.
                \end{itemize}
            \end{enumerate}
        \end{column}
        \begin{column}{0.35\textwidth}
            \textbf{Why not upload it?}\\[2pt]
            {\footnotesize Its folder's name starts with a dot. macOS hides it.\par}
        \end{column}
    \end{columns}
\end{frame}

\begin{frame}[fragile]{collect.yml: when it runs}
    \begin{columns}[T]
        \begin{column}{0.37\textwidth}
            A \textbf{workflow} is a file of steps. GitHub runs it on its own computers when something happens.

            \medskip
            Here, two things: a \texttt{schedule}, and a button, \texttt{workflow\_dispatch}. When will yours run first?

            \medskip
            {\scriptsize\textbf{Docs:} \href{https://docs.github.com/en/actions/reference/workflows-and-actions/events-that-trigger-workflows\#schedule}{schedule}\par}
        \end{column}
        \begin{column}{0.6\textwidth}
\begin{lstlisting}[language={}]
# The name on the Actions tab
name: Collect Boulder's hourly forecast

# When it runs:
on:
  # at minute 17 of every hour, in UTC
  schedule:
    - cron: "17 * * * *"
  # and when you click "Run workflow"
  workflow_dispatch:
\end{lstlisting}
        \end{column}
    \end{columns}
\end{frame}

\begin{frame}{Five fields, in UTC}
    \begin{columns}[T]
        \begin{column}{0.37\textwidth}
            \texttt{cron} is the old Unix scheduler. Actions reads its five fields, in UTC unless told otherwise. A \texttt{*} means every value.

            \medskip
            Runs at minute 0 are delayed most often. Which Boulder hour is 13:30 UTC?

            \medskip
            {\scriptsize\textbf{Try one:} \href{https://crontab.guru}{crontab.guru}\par}
        \end{column}
        \begin{column}{0.6\textwidth}
            {\footnotesize
            \begin{tabular}{@{}cccccl@{}}
                \toprule
                min & hour & day & month & weekday & \textbf{runs at} \\
                \midrule
                \texttt{17} & \texttt{*} & \texttt{*} & \texttt{*} & \texttt{*} & minute 17 of every hour \\
                \texttt{0} & \texttt{*/6} & \texttt{*} & \texttt{*} & \texttt{*} & every six hours \\
                \texttt{30} & \texttt{13} & \texttt{*} & \texttt{*} & \texttt{1-5} & 13:30, Monday to Friday \\
                \bottomrule
            \end{tabular}\par}
        \end{column}
    \end{columns}
\end{frame}

\begin{frame}[fragile]{collect.yml: permission to commit}
    \begin{columns}[T]
        \begin{column}{0.37\textwidth}
            \texttt{contents: write} lets the workflow commit to your repository.

            \medskip
            Why does it need that? What else could it do with it?

            \medskip
            {\scriptsize\textbf{Docs:} \href{https://docs.github.com/en/actions/reference/workflows-and-actions/workflow-syntax\#permissions}{permissions}\par}
        \end{column}
        \begin{column}{0.6\textwidth}
\begin{lstlisting}[language={}]
# Let the workflow commit to this repository
permissions:
  contents: write

# One job, on a fresh Ubuntu machine
jobs:
  collect:
    runs-on: ubuntu-latest
    # Its steps: the next three frames
    steps:
      ...
\end{lstlisting}
        \end{column}
    \end{columns}
\end{frame}

\begin{frame}[fragile]{collect.yml: a fresh computer each time}
    \begin{columns}[T]
        \begin{column}{0.37\textwidth}
            Each run starts on an empty Ubuntu machine. First it copies your repository's files. Then it installs Python.

            \medskip
            What would be missing if a step were left out?
        \end{column}
        \begin{column}{0.6\textwidth}
\begin{lstlisting}[language={}]
    steps:
      # Copy the repository onto the empty machine
      - name: Get this repository's files
        uses: actions/checkout@v7

      # Install Python, version 3.12
      - name: Set up Python
        uses: actions/setup-python@v7
        with:
          python-version: "3.12"
\end{lstlisting}
        \end{column}
    \end{columns}
\end{frame}

\begin{frame}[fragile]{collect.yml: run the collector}
    \begin{columns}[T]
        \begin{column}{0.37\textwidth}
            These two steps are what you would type in a terminal.

            \medskip
            Step names show in the run's log. Which step will you open when a run fails?
        \end{column}
        \begin{column}{0.6\textwidth}
\begin{lstlisting}[language={}]
      # Install the one library collector.py needs
      - name: Install requests
        run: pip install requests

      # Then run our script
      - name: Collect the forecast
        run: python collector.py
\end{lstlisting}
        \end{column}
    \end{columns}
\end{frame}

\begin{frame}[fragile]{collect.yml: commit the new rows}
    \begin{columns}[T]
        \begin{column}{0.37\textwidth}
            The last step commits \texttt{forecasts.csv} as \texttt{github-actions[bot]}, if any rows changed.

            \medskip
            If the collector fails, this step is skipped. What does your file look like then?
        \end{column}
        \begin{column}{0.6\textwidth}
\begin{lstlisting}[language={}]
      # Commit the CSV, if any rows changed
      - name: Commit the new rows
        run: |
          # Sign the commit as GitHub's bot
          git config user.name "github-actions[bot]"
          git config user.email "..."
          # Stage the CSV; commit only if it changed
          git add forecasts.csv
          git diff --staged --quiet || git commit -m "..."
          # Take any newer commits, then push
          git pull --rebase
          git push
\end{lstlisting}
        \end{column}
    \end{columns}
\end{frame}

% img/run-workflow.png: a placeholder until someone captures the real page (img/IMAGES.md).
\begin{frame}{Step 5: start the first run}
    \begin{columns}[T]
        \begin{column}{0.37\textwidth}
            A scheduled run can start late. Start the first one by hand.
            \begin{enumerate}\small\setcounter{enumi}{5}
                \item Click the \textbf{Actions} tab. If GitHub asks, enable workflows. Click \textbf{Collect Boulder's hourly forecast}. Click \textbf{Run workflow}, then the green \textbf{Run workflow}.
            \end{enumerate}
        \end{column}
        \begin{column}{0.6\textwidth}
            \centering
            \includegraphics[width=\textwidth]{img/run-workflow.png}
        \end{column}
    \end{columns}
\end{frame}

\begin{frame}{Step 5: read the log}
    \begin{columns}[T]
        \begin{column}{0.6\textwidth}
            \begin{enumerate}\small\setcounter{enumi}{6}
                \item Refresh the page until the run shows. Click the run, then the job \texttt{collect}. Click \textbf{Collect the forecast} to open that step.
            \end{enumerate}
            \medskip
            How many hours did it save? Which \texttt{updateTime}?
        \end{column}
        \begin{column}{0.35\textwidth}
            \textbf{A green check}\\[2pt]
            {\footnotesize Every step worked.\par}
            \medskip
            {\scriptsize\textbf{Docs:} \href{https://docs.github.com/en/actions/how-tos/monitor-workflows/use-workflow-run-logs}{workflow run logs}\par}
        \end{column}
    \end{columns}
\end{frame}

\begin{frame}{Step 5: see the commit}
    \begin{columns}[T]
        \begin{column}{0.6\textwidth}
            \begin{enumerate}\small\setcounter{enumi}{7}
                \item Click the \textbf{Code} tab. Find \texttt{forecasts.csv}, and the commit by \texttt{github-actions[bot]}. Click the file.
            \end{enumerate}
            \medskip
            From now on it runs by itself, every hour, until December 11.
        \end{column}
        \begin{column}{0.35\textwidth}
            \textbf{Git scraping}\\[2pt]
            {\footnotesize Each run kept as a commit, as newsrooms do. What can the history show that the file can't?\par}
        \end{column}
    \end{columns}
\end{frame}

\begin{frame}[fragile]{Step 6: your raw address}
    \begin{columns}[T]
        \begin{column}{0.37\textwidth}
            GitHub serves each file of a public repository at a \textbf{raw} address: the file alone, without GitHub's page.

            \medskip
            Put your user name in place of \texttt{<you>}.
        \end{column}
        \begin{column}{0.6\textwidth}
\begin{lstlisting}
# Your file's raw address, with your user name
my_series = (
    "https://raw.githubusercontent.com/"
    "<you>/boulder-forecast/main/forecasts.csv"
)
\end{lstlisting}
        \end{column}
    \end{columns}
\end{frame}

\begin{frame}[fragile]{Step 6: read what your collector saved}
    \begin{columns}[T]
        \begin{column}{0.37\textwidth}
            \texttt{pd.read\_csv()} reads from an address too. The raw copy can be five minutes behind.

            \medskip
            Which columns would differ from your neighbor's file?
        \end{column}
        \begin{column}{0.6\textwidth}
\begin{lstlisting}
# Read the CSV straight from its address
my_forecasts = pd.read_csv(my_series)
# How many rows and columns
print(my_forecasts.shape)
# Its first rows
my_forecasts.head()
\end{lstlisting}
        \end{column}
    \end{columns}
\end{frame}

\begin{frame}{Before Friday}
    \begin{columns}[T]
        \begin{column}{0.6\textwidth}
            Open your \textbf{Actions} tab once. Each run is a line, with a green check or a red X.

            \medskip
            How many runs do you expect by Friday's class?
        \end{column}
        \begin{column}{0.35\textwidth}
            \textbf{Friday}\\[2pt]
            {\footnotesize Part 2 starts at Step 7, with its own setup. Today's cells don't need to have run.\par}
        \end{column}
    \end{columns}
\end{frame}

% =========================================================================
\section{Friday · Live-Coding Lab}
% =========================================================================

\begin{frame}{Today: Part 2, what the collector gathered}
    \begin{columns}[T]
        \begin{column}{0.6\textwidth}
            The course has run the same \texttt{collector.py} every hour since November 13.

            \medskip
            Today we count how often the forecast changed and follow one hour as it came closer. Then we break a run on purpose.
        \end{column}
        \begin{column}{0.35\textwidth}
            \textbf{Start at Step 7}\\[2pt]
            {\footnotesize Download the new \texttt{lab-14.ipynb}. Missed Wednesday? Start here.\par}
        \end{column}
    \end{columns}
\end{frame}

\begin{frame}[fragile]{Step 7: set up again}
    \begin{columns}[T]
        \begin{column}{0.37\textwidth}
            This cell is given. Put your email in \texttt{HEADERS}, and run it.

            \medskip
            \texttt{matplotlib} is new today. Step 9 draws a plot with it.
        \end{column}
        \begin{column}{0.6\textwidth}
\begin{lstlisting}
# The libraries for Part 2
import requests
import pandas as pd
import matplotlib.pyplot as plt
# Name our program, and give your email
HEADERS = {"User-Agent":
           "WebDataScience/1.0 (your-email@colorado.edu)"}
\end{lstlisting}
        \end{column}
    \end{columns}
\end{frame}

\begin{frame}[fragile]{Step 7: two addresses}
    \begin{columns}[T]
        \begin{column}{0.37\textwidth}
            The same cell names two addresses.

            \medskip
            \texttt{series} is the course's file. A new run adds to it at minute 17 of every hour.
        \end{column}
        \begin{column}{0.6\textwidth}
\begin{lstlisting}
# The NWS's API
nws_base = "https://api.weather.gov"
# The course's series, collected every hour
series = (
    "https://raw.githubusercontent.com/"
    "cuinfoscience/INFO4617-Fall2026/main/"
    "labs/week-14/data/forecasts.csv"
)
\end{lstlisting}
        \end{column}
    \end{columns}
\end{frame}

\begin{frame}[fragile]{Step 7: the course's series}
    \begin{columns}[T]
        \begin{column}{0.37\textwidth}
            Read it, and look at its last rows.

            \medskip
            What can this file tell you that a live request can't? What can't it tell you?
        \end{column}
        \begin{column}{0.6\textwidth}
\begin{lstlisting}
# Read the course's CSV
forecasts = pd.read_csv(series)
# How many rows and columns
print(forecasts.shape)
# Its last rows
forecasts.tail()
\end{lstlisting}
        \end{column}
    \end{columns}
\end{frame}

\begin{frame}[fragile]{Step 8: runs, and forecasts}
    \begin{columns}[T]
        \begin{column}{0.37\textwidth}
            Each run is one \texttt{collected\_at}. Each forecast the NWS issued is one \texttt{update\_time}.

            \medskip
            If the forecast changed every hour, the two counts would match. Do they?
        \end{column}
        \begin{column}{0.6\textwidth}
\begin{lstlisting}
# How many runs: distinct collected_at times
print(forecasts["collected_at"].nunique(), "runs")
# How many forecasts: distinct update_time values
print(forecasts["update_time"].nunique(), "forecasts")
\end{lstlisting}
        \end{column}
    \end{columns}
\end{frame}

\begin{frame}[fragile]{Step 8: each forecast, and its runs}
    \begin{columns}[T]
        \begin{column}{0.37\textwidth}
            This cell is given. For each forecast, it counts the runs that collected it, and finds the first and last.

            \medskip
            Which forecasts lasted longest? What does that say about when forecasters update it?

            \medskip
            The notebook groups and counts in one line.
        \end{column}
        \begin{column}{0.6\textwidth}
\begin{lstlisting}
# The run times, grouped by forecast
runs = forecasts.groupby("update_time")["collected_at"]
# For each: how many runs, the first, and the last
per_forecast = runs.agg([
    "nunique",
    "min",
    "max",
])
# Print every row of the table
print(per_forecast.to_string())
\end{lstlisting}
        \end{column}
    \end{columns}
\end{frame}

\begin{frame}[fragile]{Microsprint: the share that was new}
    \begin{columns}[T]
        \begin{column}{0.37\textwidth}
            Divide forecasts by runs: the share of runs that brought a forecast no earlier run had.

            \medskip
            Thirty collectors each ask 24 times a day. How many of their 720 requests bring something new?

            \medskip
            The notebook divides in one line.
        \end{column}
        \begin{column}{0.6\textwidth}
\begin{lstlisting}
# How many forecasts, and how many runs
n_forecasts = forecasts["update_time"].nunique()
n_runs = forecasts["collected_at"].nunique()
# The share of runs that brought a new forecast
new_share = n_forecasts / n_runs
print(round(new_share, 2))
# The same share of 720 requests a day
print(round(720 * new_share), "of 720 daily requests")
\end{lstlisting}
        \end{column}
    \end{columns}
\end{frame}

\begin{frame}{What a bot costs a server}
    \begin{columns}[T]
        \begin{column}{0.6\textwidth}
            Step 8 counted how often the forecast changed. Our collector asks 24 times a day anyway. On November 12, an hourly collector got the same forecast 12 times in a row.

            \medskip
            \texttt{Last-Modified} and \texttt{Cache-Control} say how new an answer is. Who pays for the requests that bring nothing new?
        \end{column}
        \begin{column}{0.35\textwidth}
            \textbf{Ask less?}\\[2pt]
            {\footnotesize What would a collector need to know to ask less often and miss nothing?\par}
        \end{column}
    \end{columns}
\end{frame}

\begin{frame}[fragile]{Step 9: one hour, forecast again and again}
    \begin{columns}[T]
        \begin{column}{0.37\textwidth}
            Every run forecasts the next six and a half days. So each hour is forecast by many runs.

            \medskip
            Take noon on Wednesday, our first class. How many runs forecast it?

            \medskip
            The notebook picks these rows in one line.
        \end{column}
        \begin{column}{0.6\textwidth}
\begin{lstlisting}
# Noon on Wednesday, November 18, in Boulder's time
noon = "2026-11-18T12:00:00-07:00"
# True for each row that forecasts that hour
is_noon = forecasts["start_time"] == noon
# Keep those rows, and count them
noon_rows = forecasts[is_noon]
print(len(noon_rows))
\end{lstlisting}
        \end{column}
    \end{columns}
\end{frame}

\begin{frame}[fragile]{Step 9: lead time}
    \begin{columns}[T]
        \begin{column}{0.37\textwidth}
            A row's \textbf{lead time} is how far ahead its run was: the hour's \texttt{start\_time} minus the run's \texttt{collected\_at}.

            \medskip
            This cell is given. The notebook subtracts in one line, without naming \texttt{hour} and \texttt{run}.
        \end{column}
        \begin{column}{0.6\textwidth}
\begin{lstlisting}
# Noon, and when each run asked, as times in UTC
hour = pd.to_datetime(noon_rows["start_time"], utc=True)
run = pd.to_datetime(noon_rows["collected_at"], utc=True)
# Lead time: how far ahead each run was
lead = hour - run
# The same, in hours
lead_hours = lead.dt.total_seconds() / 3600
\end{lstlisting}
        \end{column}
    \end{columns}
\end{frame}

\begin{frame}[fragile]{Step 9: plot it}
    \begin{columns}[T]
        \begin{column}{0.37\textwidth}
            The same given cell plots the forecast temperature against the lead time.

            \medskip
            Where does the line settle?
        \end{column}
        \begin{column}{0.6\textwidth}
\begin{lstlisting}
# One figure, wide and short
fig, ax = plt.subplots(figsize=(8, 4))
# Temperature against lead time, a dot per run
ax.plot(
    lead_hours,
    noon_rows["temperature"],
    marker=".",
)
# Furthest ahead on the left, nearest on the right
ax.invert_xaxis()
# Then a title, axis labels, and plt.show()
...
\end{lstlisting}
        \end{column}
    \end{columns}
\end{frame}

\begin{frame}[fragile]{Step 9: furthest ahead, and nearest}
    \begin{columns}[T]
        \begin{column}{0.37\textwidth}
            Print the first row and the last: the furthest ahead and the nearest.

            \medskip
            Which one was right? What else would you have to collect to know?

            \medskip
            The notebook writes the list of columns into each line.
        \end{column}
        \begin{column}{0.6\textwidth}
\begin{lstlisting}
# The three columns to show
columns = [
    "collected_at",
    "temperature",
    "short_forecast",
]
# The first run, furthest ahead
print(noon_rows[columns].head(1))
# The last run, nearest
print(noon_rows[columns].tail(1))
\end{lstlisting}
        \end{column}
    \end{columns}
\end{frame}

\begin{frame}[fragile]{Microsprint: midnight}
    \begin{columns}[T]
        \begin{column}{0.37\textwidth}
            Do the same for midnight that night.

            \medskip
            How far apart were its lowest and highest forecasts? Is that more or less than noon's?

            \medskip
            The notebook picks the rows in one line.
        \end{column}
        \begin{column}{0.6\textwidth}
\begin{lstlisting}
# Midnight's rows, picked as for noon, and how many
midnight = "2026-11-19T00:00:00-07:00"
is_midnight = forecasts["start_time"] == midnight
midnight_rows = forecasts[is_midnight]
print(len(midnight_rows))
# Its lowest and highest forecast temperature
print(
    midnight_rows["temperature"].min(),
    midnight_rows["temperature"].max(),
)
\end{lstlisting}
        \end{column}
    \end{columns}
\end{frame}

\begin{frame}[fragile]{Step 10: is it still running?}
    \begin{columns}[T]
        \begin{column}{0.37\textwidth}
            A collector that stops doesn't say so. Its file just stops growing.

            \medskip
            A \textbf{freshness check} asks how old the newest row is. Under about an hour and a half is fresh. Is the course's?
        \end{column}
        \begin{column}{0.6\textwidth}
\begin{lstlisting}
# The newest run's time
newest = pd.to_datetime(forecasts["collected_at"]).max()
# That time, and how long ago it was
print(newest, pd.Timestamp.now(tz="UTC") - newest)
\end{lstlisting}
        \end{column}
    \end{columns}
\end{frame}

\begin{frame}[fragile]{Step 10: is yours?}
    \begin{columns}[T]
        \begin{column}{0.37\textwidth}
            Check your own file the same way, from Wednesday's raw address.

            \medskip
            Is it fresh? If not, where would you look first?

            \medskip
            The notebook reads the file and finds its newest run in one line.
        \end{column}
        \begin{column}{0.6\textwidth}
\begin{lstlisting}
# Your raw address, as on Wednesday
my_series = (
    "https://raw.githubusercontent.com/"
    "<you>/boulder-forecast/main/forecasts.csv"
)
# Read your file
my_rows = pd.read_csv(my_series)
# Its newest run, and how long ago it was
my_newest = pd.to_datetime(my_rows["collected_at"]).max()
print(my_newest, pd.Timestamp.now(tz="UTC") - my_newest)
\end{lstlisting}
        \end{column}
    \end{columns}
\end{frame}

\begin{frame}[fragile]{Step 10: is the forecast fresh?}
    \begin{columns}[T]
        \begin{column}{0.37\textwidth}
            Rows can keep coming while \texttt{update\_time} stands still. Half a day without a change can be normal.

            \medskip
            How long ago did the forecasters last change it? How old would you call stale? Where could this check run to warn you?

            \medskip
            The notebook prints this in one line.
        \end{column}
        \begin{column}{0.6\textwidth}
\begin{lstlisting}
# When the forecasters last changed the forecast
updated = pd.to_datetime(forecasts["update_time"]).max()
# How long ago that was
print(pd.Timestamp.now(tz="UTC") - updated)
\end{lstlisting}
        \end{column}
    \end{columns}
\end{frame}

\begin{frame}[fragile]{Step 11: a mistyped address}
    \begin{columns}[T]
        \begin{column}{0.37\textwidth}
            This request is \textbf{designed to break}. A deliberate error needs no GitHub issue; only code that is broken by accident does.

            \medskip
            Where is the typo?
        \end{column}
        \begin{column}{0.6\textwidth}
\begin{lstlisting}
# The hourly forecast's address, mistyped
broken_url = (
    f"{nws_base}/gridpoints/BOU/54,75/"
    "forecast/hourlyy"
)
# Send it
broken_response = requests.get(
    broken_url,
    headers=HEADERS,
)
\end{lstlisting}
        \end{column}
    \end{columns}
\end{frame}

\begin{frame}[fragile]{Step 11: what came back}
    \begin{columns}[T]
        \begin{column}{0.37\textwidth}
            Read the answer, and print its status code and \texttt{detail}.

            \medskip
            What status code? Does the \texttt{detail} say what was wrong?
        \end{column}
        \begin{column}{0.6\textwidth}
\begin{lstlisting}
# The answer, as a dictionary
broken = broken_response.json()
# The status code, and the NWS's error
print(broken_response.status_code, broken.get("detail"))
\end{lstlisting}
        \end{column}
    \end{columns}
\end{frame}

% img/failed-run.png: a placeholder until someone captures the real page (img/IMAGES.md).
\begin{frame}{Step 11: break your collector}
    \begin{columns}[T]
        \begin{column}{0.37\textwidth}
            \begin{enumerate}\small
                \item Edit \texttt{collector.py}. In \texttt{URL}, change \texttt{hourly} to \texttt{hourlyy}. Commit the change.
                \item Run the workflow. Open the run with the red X, then its job.
            \end{enumerate}
            \medskip
            Which line of the log names the error? Was the commit step run?
        \end{column}
        \begin{column}{0.6\textwidth}
            \centering
            \includegraphics[width=\textwidth]{img/failed-run.png}
        \end{column}
    \end{columns}
\end{frame}

\begin{frame}{Step 11: fix it}
    \begin{columns}[T]
        \begin{column}{0.6\textwidth}
            \begin{enumerate}\small\setcounter{enumi}{2}
                \item Change \texttt{hourlyy} back to \texttt{hourly}. Commit the change. Run the workflow again, and wait for the green check.
            \end{enumerate}
            \medskip
            A red X is a loud failure: GitHub can email whoever made the workflow. Did you get an email?

            \medskip
            {\scriptsize\textbf{Docs:} \href{https://docs.github.com/en/actions/concepts/workflows-and-actions/notifications-for-workflow-runs}{notifications for workflow runs}\par}
        \end{column}
        \begin{column}{0.35\textwidth}
        \end{column}
    \end{columns}
\end{frame}

\begin{frame}[fragile]{Microsprint: not every mistype fails}
    \begin{columns}[T]
        \begin{column}{0.37\textwidth}
            Ask for the hourly forecast with the grid numbers swapped, \texttt{75,54}.

            \medskip
            Do you expect an error?
        \end{column}
        \begin{column}{0.6\textwidth}
\begin{lstlisting}
# The grid numbers swapped
swapped_grid_url = (
    f"{nws_base}/gridpoints/BOU/75,54/"
    "forecast/hourly"
)
# Send it
swapped_grid_response = requests.get(
    swapped_grid_url,
    headers=HEADERS,
)
\end{lstlisting}
        \end{column}
    \end{columns}
\end{frame}

\begin{frame}[fragile]{Microsprint: what came back}
    \begin{columns}[T]
        \begin{column}{0.37\textwidth}
            Print the status code and the first hour's temperature.

            \medskip
            What status code? Did anything tell you it wasn't Boulder?

            \medskip
            The notebook prints the temperature without naming it first.
        \end{column}
        \begin{column}{0.6\textwidth}
\begin{lstlisting}
# The part that holds the forecast
swapped_grid = swapped_grid_response.json()["properties"]
# The first hour's temperature
temperature = swapped_grid["periods"][0]["temperature"]
# The status code, and that temperature
print(swapped_grid_response.status_code, temperature)
\end{lstlisting}
        \end{column}
    \end{columns}
\end{frame}

\begin{frame}{Quiet failures}
    \begin{columns}[T]
        \begin{column}{0.6\textwidth}
            Some failures leave no red X. A square near Aurora answers 200. A forecast can stop changing. GitHub pauses a public repository's schedule after 60 days without activity.

            \medskip
            Which would you rather find a month from now: a gap, or a month of Aurora's forecasts?
        \end{column}
        \begin{column}{0.35\textwidth}
            \centering
            \includegraphics[width=.9\textwidth]{slides/week-02/img/this-is-fine.jpg}
        \end{column}
    \end{columns}
\end{frame}

% Free tiers checked on 2026-11-13 on each provider's own pages; sources in img/IMAGES.md ("Sources for
% the hosted-scheduler table").
\begin{frame}{Other places to run a schedule}
    Actions is one of several hosted schedulers. Each free tier has its own limit. Which of these could run our collector every hour, for free?

    \medskip
    {\footnotesize\renewcommand{\arraystretch}{1.25}
    \begin{tabular}{@{}lll@{}}
        \toprule
        \textbf{Service} & \textbf{What runs} & \textbf{Free tier, read on Nov 13, 2026} \\
        \midrule
        GitHub Actions (ours) & A workflow's steps & Free in public repositories; every 5 minutes at most \\
        GitLab scheduled pipelines & A CI/CD pipeline & 400 compute minutes a month; 24 runs a day \\
        PythonAnywhere scheduled tasks & A Python script & None on free accounts made since January 2026 \\
        Cloudflare Workers Cron Triggers & A Worker, in JavaScript or Python & 5 cron triggers; 10 ms of CPU a run \\
        Google Cloud Scheduler & A call to a URL you give it & 3 jobs a month per billing account \\
        Val Town & A val, in TypeScript & Every 15 minutes at most \\
        \bottomrule
    \end{tabular}\par}
\end{frame}

\begin{frame}[fragile]{cron on your own computer}
    \begin{columns}[T]
        \begin{column}{0.37\textwidth}
            \texttt{crontab -e} opens your computer's schedule. This line runs the collector at minute 17, in your computer's time zone.

            \medskip
            It runs only while the computer is awake. What happens when your laptop sleeps?
        \end{column}
        \begin{column}{0.6\textwidth}
\begin{lstlisting}[language={},literate={~}{{\raisebox{-0.6ex}{\textasciitilde}}}1]
# At minute 17 of every hour, run the collector
17 * * * * cd ~/boulder-forecast && python3 collector.py
\end{lstlisting}
            \medskip
            {\footnotesize Wrong Python? Give the full path that \texttt{which python3} prints in \texttt{webdata}.\par}
        \end{column}
    \end{columns}
\end{frame}

\begin{frame}{Step 12: what did we learn?}
    \begin{columns}[T]
        \begin{column}{0.6\textwidth}
            Answer in the notebook, as comments.
            \begin{itemize}\small
                \item How often should a collector ask?
                \item How far did noon's forecast move?
                \item Name a failure with a red X, and one without.
                \item What should your README say about when you collected?
            \end{itemize}
        \end{column}
        \begin{column}{0.35\textwidth}
            \textbf{Due Sunday, Nov 22}\\[2pt]
            {\footnotesize Chapter 14's Recommended Exercises, from its companion notebook. Graded on a good-faith attempt.\par}
        \end{column}
    \end{columns}
\end{frame}

% Checkpoint 4's criteria: canvas/course_spec.py (PORTFOLIO_WEEKS and _PORTFOLIO_CRITERIA, week 14);
% due date: docs/course-facts.md.
\begin{frame}{Friday's last 15 minutes: portfolio checkpoint 4}
    \dayband{Due Sunday, Nov 29}{Draft your write-up}

    \medskip
    \begin{itemize}\small
        \item \textbf{A full draft} of \texttt{report.md} or \texttt{report.pdf}: every section, your own results and figures, your sources cited
        \item \textbf{All four files}, present and current: \texttt{README.md}, \texttt{proposal.md}, your notebook, and your report
        \item \textbf{A README} that lets someone else reproduce the project: setup, how to run it, where the data came from
    \end{itemize}

    \medskip
    Work on your README's \textbf{How to reproduce} now: when you sent your requests, and from what code. Next: final presentations, after Fall Break.
\end{frame}

% =========================================================================
\begin{frame}{Key takeaways}
    \begin{enumerate}
        \item One request is a \textbf{snapshot}. A series needs a program on a schedule.
        \item A collector \textbf{appends}, and records when it asked.
        \item \textbf{Count} what changed before deciding how often to ask.
        \item A red X is a \textbf{loud} failure. Check freshness for the quiet ones.
        \item Write a \textbf{last day} into every collector.
    \end{enumerate}
\end{frame}

\end{document}
